\documentclass[sigconf,preprint]{acmart}

\usepackage[capitalise]{cleveref}

% TODO before submission:
%  - replace [preprint] with [anonymous,review] for double-blind venues;
%  - fill in the real \acmConference/\acmISBN/\acmDOI from the rights form;
%  - re-generate and verify CCS codes at https://dl.acm.org/ccs.
\settopmatter{printacmref=false}
\renewcommand\footnotetextcopyrightpermission[1]{}
\acmConference[Preprint]{Preprint --- venue TBD}{2026}{}
\acmYear{2026}
\acmISBN{}
\acmDOI{}
\setcopyright{none}

\title{AlignmentGraph-ISD: An AI Agent Harness for Controlled Continuous Refinement of LLM-Generated Training Programs}

% TODO: confirm supervisor emails and exact affiliation wording before submission.
\author{Phuong Viet Bui}
\affiliation{%
  \department{School of Information and Communication Technology}
  \institution{Hanoi University of Science and Technology}
  \city{Hanoi}
  \country{Vietnam}
}
\email{phuong.bv242011m@sis.hust.edu.vn}

\author{Hai Van Pham}
\authornote{Primary supervisor.}
\affiliation{%
  \department{School of Information and Communication Technology}
  \institution{Hanoi University of Science and Technology}
  \city{Hanoi}
  \country{Vietnam}
}
\email{haipv@soict.hust.edu.vn}

\author{Long Kim Cu}
\authornote{Co-supervisor.}
\affiliation{%
  \institution{Center for Information Technology, Ministry of Science and Technology}
  \city{Hanoi}
  \country{Vietnam}
}

\renewcommand{\shortauthors}{Bui et al.}

\begin{document}

\begin{abstract}
Current commercial LLMs already provide strong support for instructional systems
design (ISD). However, dependence on third-party APIs, high costs, and data
safety risks remain barriers to practical deployment in educational institutions.
Open-weight SLMs deployed on-premise offer a more viable path, but as models
become smaller, the quality burden shifts from the model to the system around
it---above all constructive alignment, the coherence among objectives,
activities, and assessments, which existing agent systems leave for the model to
maintain implicitly. Research on ISD agents has not yet
treated quality retention under model downscaling as a central problem. We therefore
propose AlignmentGraph-ISD, a multi-agent system (MAS) that organizes the ISD process
around an alignment graph: explicit, checkable alignment state that takes this
coherence burden off the model. We evaluate
it with ISD-Agent-Bench and a deterministic, no-LLM alignment-score measurement
across a four-size Qwen model ladder (0.8B--9B). As models shrink, the MAS remains competitive and
robust, with a 0\% failure rate, satisfies the panel advantage
rule on alignment scores ($\geq$4 of 6 directional scores are
significantly positive at every size), and leads on ADDIE score at the three
smaller sizes (\rqOneHarnessAddieTwoB{} versus \rqOneBestBaselineAddieTwoB{} for
the strongest baseline at 2B). Its gap widens significantly toward smaller models
before reaching parity with the strongest baseline at 9B. This robustness comes
with significantly lower Trajectory scores than the leading single-loop systems
and token costs above the baseline median. Beyond quantitative measurement, the MAS's explicit
alignment graph lets instructional designers audit and refine final outputs,
supporting teacher-in-the-loop practice.
\end{abstract}


\begin{CCSXML}
<ccs2012>
   <concept>
       <concept_id>10010405.10010489</concept_id>
       <concept_desc>Applied computing~Education</concept_desc>
       <concept_significance>500</concept_significance>
   </concept>
   <concept>
       <concept_id>10010147.10010178.10010219.10010220</concept_id>
       <concept_desc>Computing methodologies~Multi-agent systems</concept_desc>
       <concept_significance>500</concept_significance>
   </concept>
   <concept>
       <concept_id>10010147.10010178.10010179</concept_id>
       <concept_desc>Computing methodologies~Natural language processing</concept_desc>
       <concept_significance>300</concept_significance>
   </concept>
</ccs2012>
\end{CCSXML}

\ccsdesc[500]{Applied computing~Education}
\ccsdesc[500]{Computing methodologies~Multi-agent systems}
\ccsdesc[300]{Computing methodologies~Natural language processing}

\keywords{instructional systems design, LLM agents, agent harness, constructive alignment, multi-agent orchestration, multi-session refinement}

\maketitle

\section{Introduction}
\label{sec:introduction}

Commercial large language models (LLMs) can produce high-quality curricula and
support instructional systems design (ISD) from needs analysis through assessment
~\cite{isdagentbench}. Institutional use nevertheless brings recurring inference
cost, external-API dependence, and third-party transfer of learner or institutional
data. On-premise open-weight small language models (SLMs) offer more control, but as
the model shrinks, the harness that manages state, coordinates calls, and checks
outputs must carry more of the system-level burden.

Existing ISD and lesson-planning agents emphasize automation, personalization, or
quality with capable foundation models. They do not treat preservation of
instructional-design quality during model downscaling as a joint design and
evaluation objective. A deployment-oriented study must establish both quality at
each feasible size and whether the relative gap changes as capacity falls.

A second gap concerns constructive alignment among learning objectives, activities,
and assessments~\cite{biggs1996enhancing,biggs2011teaching}. ISD benchmarks assess
this property mainly through LLM-judge rubrics~\cite{isdagentbench}, placing the
measurement instrument in the same broad family as the generator. Porter/Webb
indices and embedding methods can quantify correspondence to human-authored
standards~\cite{porter2002measuring,webb1997criteria}, but no ISD benchmark uses
such methods to examine the intrinsic alignment of a wholly machine-generated
design, where both sides of every relation are system outputs.

Our response begins from the deployment constraint. Each model call receives one
bounded ISD concern rather than the full open-ended, multi-constraint task. The
premise that narrower subproblems can be more tractable is inherited from work on
decomposed prompting and small-model orchestration
~\cite{zhou2023leasttomost,khot2023decomposed,ferraz2024decrim,
shen2024multillmagent,belcak2025slm}; it is not tested here because every ablation
retains decomposition. Nor do we assume that its benefit grows as models shrink:
harness effects can be non-monotonic across model tiers~\cite{cho2026harness}, so
the scale interaction is an empirical question.

Decomposition creates its own problem: separately generated artifacts must retain
one pedagogical intent. AlignmentGraph-ISD therefore uses constructive alignment to
shape a shared graph of objectives, activities, assessments, and their relations.
Specialists read only relevant graph context and update scoped state; a verifier
checks the evolving design, and a graph dump exposes the final relations for audit.
The graph is a substrate for constrained generation, not a claim that prose is
aligned by construction (\Cref{sec:method}).

The evaluation separates quality level from scale interaction. RQ1a compares the
harness with six agents at each Qwen size; RQ1b tests whether those gaps widen as
the model shrinks. The \texttt{addie\_median}, \texttt{trajectory\_score}, and
\texttt{total\_score} judge signals receive identical analyses and equal reporting
status. ADDIE leads the narrative only because its artifact rubric matches RQ1's
content-quality construct. RQ2 asks how measurable signals of constructive
alignment vary across the same ladder, without claiming that automated scores
validate the construct itself (\Cref{sec:rqs,sec:protocol}).

This paper makes two contributions:
\begin{enumerate}
  \item \textbf{AlignmentGraph-ISD}, a multi-agent ADDIE harness whose alignment
  graph is explicit, checkable state and an auditable educator-facing artifact.
  \item \textbf{An automated constructive-alignment evaluation protocol} extending
  ISD-Agent-Bench with a no-LLM, non-hierarchical panel of seven signals from
  mechanically independent textual-correspondence and cognitive-demand instrument
  families. It operationalizes established criteria rather than proposing a new
  construct~\cite{biggs2011teaching,qm2018rubric,porter2002measuring,
  webb1997criteria}; cross-family agreement is an initial validity bridge, while
  expert-human validation remains future work.
\end{enumerate}

The experiment spans four Qwen sizes from 0.8B to 9B, three independent runs per
cell, and three ablations inside the ladder, with between-run variation reported
for every estimate. Following this Introduction, \Cref{sec:related-work} locates
the two gaps; \Cref{sec:method} presents the harness and evaluation protocol;
\Cref{sec:experimental-methodology,sec:results} describe the experimental design
and findings; \Cref{sec:discussion} interprets the results and bounds their
validity; and \Cref{sec:conclusion} closes the paper.

\section{Related Work}
\label{sec:related-work}

\subsection{LLM Agents for Instructional Design}

The systems closest to ours apply ISD theory to agent architecture.
ISD-Agent-Bench~\cite{isdagentbench} implements four theory-grounded agents:
ReAct-ADDIE (a ReAct loop over ISD tools~\cite{yao2023react}), ADDIE-Agent (a
linear five-phase state graph~\cite{branch2009addie}), Dick-Carey-Agent (the
systems approach with a formative-evaluation loop~\cite{dick1978systematic}), and
RPISD-Agent (rapid prototyping~\cite{tripp1990rapid}). EduPlanner~\cite{zhang2025eduplanner}
uses adversarial generate--evaluate--optimize collaboration among agents;
LessonPlanner~\cite{lessonplanner2024} and GPT-4 teaching-plan
studies~\cite{hu2024teachingplan} improve single artifacts; LearnLM~\cite{learnlm2024}
tunes the model itself for pedagogy. All of these accept a brief and optimize the
draft produced by a design pass, and all treat constructive alignment as something
to be \emph{encouraged} through prompt design or scored after the fact---none
represents alignment as state that the system checks per objective and repairs.

\subsection{Benchmarks and Evaluation}

OpenLearnLM~\cite{openlearnlm2026}, EduBench~\cite{edubench2025},
MathTutorBench~\cite{mathtutorbench2025}, and TutorBench~\cite{tutorbench2025}
evaluate educational capabilities at the level of individual tasks.
ISD-Agent-Bench~\cite{isdagentbench} goes further by scoring both artifact quality
against an ADDIE rubric and the quality of the agent's trajectory, over synthetic
scenarios spanning a Context~$\times$~ISD matrix, with multi-judge grading to
reduce self-preference bias. We adopt it as our evaluation substrate for both
research questions, and additionally report a deterministic \emph{residual
alignment-defect audit} that measures a property the rubric does not directly
score.

\subsection{Consistency Control in Multi-Artifact Generation}

Theoretically, constructive alignment~\cite{biggs1996enhancing} defines the
consistency standard among objectives, activities, and assessments, and the
systems approach~\cite{dick1978systematic} justifies systematic verification.
Bloom's taxonomy~\cite{anderson2001taxonomy} provides the cognitive-level standard
an assessment must match. Technically, verifier and self-repair mechanisms such as
Self-Refine~\cite{madaan2023selfrefine} and Reflexion~\cite{shinn2023reflexion}
demonstrate the value of checking and correcting model output from outside the
model---value that grows as base models saturate and as smaller models are
deployed~\cite{cho2026harness,belcak2025slm}. In existing ISD systems, however,
alignment is only encouraged through prompts or scored after generation; it is
never represented as first-class state---a data structure whose links the system
can check individually, whose affected regions it can mark on change, and through
which it can propagate repairs under control.

\subsection{Synthesis and Positioning}

\Cref{tab:related} summarizes the gap. Current ISD agents optimize a single pass,
are evaluated on artifact and trajectory quality, and lack an explicit
representation of alignment for control; none is evaluated on the SLM setting that
increasingly governs deployment cost. AlignmentGraph-ISD fills the gap by making
alignment the first-class graph state of the harness, decomposing generation
around it so alignment holds by construction, verifying and repairing pedagogical
constraints as per-agent tools, and evaluating the result on both an LLM and an
SLM.

\begin{table}[t]
  \caption{Positioning against representative ISD agents.
  AbC = alignment represented as \emph{checkable state} (not just prompted);
  Ver = per-objective pedagogical verification;
  Rep = automatic repair of detected violations;
  SLM = evaluated on a small language model.}
  \label{tab:related}
  \small
  \begin{tabular}{@{}lcccc@{}}
    \toprule
    Agent & AbC & Ver & Rep & SLM \\
    \midrule
    Baseline (single prompt)               & \ding{55} & \ding{55} & \ding{55} & \ding{55} \\
    ReAct-ADDIE~\cite{isdagentbench,yao2023react}       & \ding{55} & \ding{55} & \ding{55} & \ding{55} \\
    ADDIE-Agent~\cite{branch2009addie}     & \ding{55} & \ding{55} & \ding{55} & \ding{55} \\
    Dick-Carey-Agent~\cite{dick1978systematic} & \ding{55} & (\ding{51}) & \ding{55} & \ding{55} \\
    RPISD-Agent~\cite{tripp1990rapid}      & \ding{55} & \ding{55} & \ding{55} & \ding{55} \\
    EduPlanner~\cite{zhang2025eduplanner}  & \ding{55} & (\ding{51}) & (\ding{51}) & \ding{55} \\
    \textbf{AlignmentGraph-ISD} (ours)     & \ding{51} & \ding{51} & \ding{51} & \ding{51} \\
    \bottomrule
  \end{tabular}
\end{table}

\section{The AlignmentGraph-ISD Harness}
\label{sec:harness}

\subsection{Overview}
\label{sec:overview}

The analysis in \cref{sec:introduction} requires supporting the full lifecycle of a training program: turning a vague need into a design brief, generating a complete design from that brief, and turning every later refinement request into a new version of the design without breaking constructive alignment. The observation that drives our design is that everything this lifecycle demands---persistent state, quality control, change scoping, and specialist coordination---is not a capability of the model but a responsibility of the execution environment around it. AlignmentGraph-ISD therefore concentrates these responsibilities in a harness, leaving the model to generate content and reason within that environment. A complete agent is the combination of the two (\emph{agent = harness + model}), and every design decision in this section belongs to the harness, independent of whether the model is an LLM or a low-cost SLM.

\begin{figure*}[t]
  \centering
  \fbox{\parbox{0.92\textwidth}{\centering\vspace{6em}
    [FIGURE PENDING: architecture overview --- elicitation $\rightarrow$ orchestrator and specialist agents $\rightarrow$ alignment graph (center) $\rightarrow$ verifier and cascade-scoped revision $\rightarrow$ four-layer memory; retrieval agents plugged in at the side.]
  \vspace{6em}}}
  \caption{Architecture of the AlignmentGraph-ISD harness. All modules read or write the versioned alignment graph at the center.}
  \label{fig:architecture}
\end{figure*}

The center of the harness is a versioned \emph{alignment graph} (\cref{sec:alignment-graph}); every other module reads or writes this state. A training program flows through the harness as shown in \cref{fig:architecture}: the designer's need is made concrete through a conversational elicitation loop (\cref{sec:elicitation}); from the approved brief, the orchestrator dispatches tasks to specialist agents that build the graph (\cref{sec:orchestration}), and when domain material is missing, the specialist itself decides to retrieve from pluggable retrieval agents (\cref{sec:agentic-rag}); every change to the graph---whether from initial generation, repair, or a refinement request---is verified and propagated under control (\cref{sec:verification}); the result of each session is committed as a version in multi-session memory (\cref{sec:memory}), which becomes the starting point of the next session. All modules are configuration-controlled; the configurations used in each experiment are stated in \cref{sec:experiments}.

\subsection{Alignment Graph}
\label{sec:alignment-graph}

The root of the refinement dilemma is that the relations among design components exist only implicitly in text: a reader may \emph{sense} that an assessment measures an objective, but the system has no way to check it, let alone answer what a change affects. Controlled refinement requires alignment to become \emph{data}---each link checkable, each affected region markable. AlignmentGraph-ISD therefore reifies constructive alignment~\cite{biggs1996enhancing} as an explicit data structure: at each step $t$, the design state is a graph $G_t = (V_t, E_t)$ whose vertex set covers the artifacts of all five ADDIE phases:
\begin{multline*}
V_t = V_{\mathrm{needs}} \cup V_{\mathrm{objective}} \cup V_{\mathrm{assessment}} \cup V_{\mathrm{activity}} \\
 \cup V_{\mathrm{material}} \cup V_{\mathrm{implementation}} \cup V_{\mathrm{evaluation}}.
\end{multline*}
The edge set is kept minimal: three many-to-many relations centered on objectives,
\begin{align*}
E_t &= E_{\mathrm{measures}} \cup E_{\mathrm{supports}} \cup E_{\mathrm{evaluates}},\\
E_{\mathrm{measures}} &\subseteq V_{\mathrm{assessment}} \times V_{\mathrm{objective}},\\
E_{\mathrm{supports}} &\subseteq V_{\mathrm{activity}} \times V_{\mathrm{objective}},\\
E_{\mathrm{evaluates}} &\subseteq V_{\mathrm{evaluation}} \times V_{\mathrm{objective}}.
\end{align*}
An objective is \emph{aligned} if and only if at least one assessment measures it, one activity supports it, and one evaluation evaluates it: for all $o \in V_{\mathrm{objective}}$,
\begin{multline*}
\mathrm{aligned}(o) \iff \big(\exists a\colon (a,o) \in E_{\mathrm{measures}}\big) \\
 \wedge \big(\exists x\colon (x,o) \in E_{\mathrm{supports}}\big) \wedge \big(\exists e\colon (e,o) \in E_{\mathrm{evaluates}}\big).
\end{multline*}

Only objectives are decomposed into element nodes per learning outcome (\texttt{objective:LO-*}), because constructive alignment must be checked per objective; other node types remain at section granularity. Node types that carry no static alignment edges (needs context, implementation plan, material) affect the design through the cascade operator,
\[
\mathrm{cascade}(n) = \{\, m \in V_t : \mathrm{type}(m) \in \mathrm{IMPACT}(\mathrm{type}(n)) \,\},
\]
where $\mathrm{IMPACT}$ is the impact policy mapping each node type to the set of downstream types that depend on it; every node in $\mathrm{cascade}(n)$ is marked \emph{stale}. Each node carries \texttt{stale} and \texttt{repair\_history} fields and each edge a \texttt{verified} field---these state fields are what make alignment measurable in \cref{sec:experiments}.

\subsection{Multi-Agent Orchestration and Dynamic Workflow}
\label{sec:orchestration}

Given explicit state, the next question is who operates on it and in what order. Two properties of the problem rule out a fixed pipeline: each ISD artifact type demands different expertise---needs analysis, rubric design, and implementation planning are distinct skills---and in a refinement round the work to be done depends on the \emph{current state} of the design, which a hard-coded ADDIE sequence cannot express (``regenerate only the stale part'' is not a step in any fixed pipeline). Coordination must therefore read the graph state to decide the next task, and each task must be assigned to the right specialist with a bounded write scope.

Concretely, each ISD concern has a specialist agent (needs context, objective design, assessment design, activity and material, implementation plan, evaluation and revision), plus a pedagogical verifier and an ADDIE aggregator. The meta-agent maintains a dynamic task graph $T_t = \{\tau_1, \dots, \tau_k\}$ with tasks $\tau = (\mathrm{type}, \mathrm{owner}, \mathrm{prio}, \mathrm{deps})$; at each step the planner policy $\pi$ selects the next feasible task from the graph state, and the owning agent executes it:
\[
\tau^* = \pi(G_t, T_t), \qquad G_{t+1} = A_{\mathrm{owner}(\tau^*)}\big(G_t,\, c_{\tau^*}\big).
\]
Each agent is a graph-update operator whose write scope is bounded by its specialty:
\[
A_i : (G_t, c) \mapsto G_{t+1}, \qquad \Delta(G_t, G_{t+1}) \subseteq V^{(i)} \cup E^{(i)}
\]
(e.g., the assessment agent writes only \texttt{assessment} nodes and \textsl{measures} edges). The planner creates tasks only for missing, stale, or invalid nodes---a complete, non-stale graph yields zero tasks---so the cost of a session scales with the size of the stale set rather than the size of the program. ADDIE serves only as the output contract for benchmark compatibility, not as an execution-order constraint.

\subsection{Pedagogical Verification and Cascade-Scoped Revision}
\label{sec:verification}

The second horn of the dilemma---local edits break alignment---stems from two deficits: alignment violations go undetected, and the ripple effects of a change go unhandled. The remedy must therefore be a closed loop---detect violations, route them to the right specialist, propagate impact, re-verify---which is, in essence, the formative evaluation--revision loop of the systems approach~\cite{dick1978systematic}, automated over the graph structure.

After every generation or repair turn, two independent check layers run: schema validation (JSON structure, required fields, types, identifiers) and pedagogical verification with seven constraint families (objective coverage, Bloom-verb usage~\cite{anderson2001taxonomy}, activity support, rubric consistency, feasibility, actionability, and evaluation alignment). Let $\Omega_t = \mathcal{V}(G_t)$ be the set of violations detected on $G_t$, each violation $\omega = (\mathrm{type}, x)$; the function $\mathrm{owner}(\cdot)$ routes a violation to its specialist. One revision turn, and the termination condition, are
\begin{gather*}
G_{t+1} = \mathrm{cascade}\Big(A_{\mathrm{owner}(\mathrm{type}(\omega))}\big(G_t,\, \omega\big)\Big), \quad \omega \in \Omega_t,\\
\text{terminating when } \Omega_t = \emptyset \ \lor\ t \geq B,
\end{gather*}
where $B$ is a retry budget: the loop either converges to a violation-free fixed point or stops in a controlled manner recorded in the trajectory. A refinement session enters the same path through the \emph{refinement entrypoint}: a refinement request $r$ (a target node plus a natural-language change request) marks the target stale, applies $\mathrm{cascade}(\cdot)$ to scope the impact, and lets the loop above regenerate exactly the stale set. We call the pair ``scope-determining invalidation + in-scope regeneration'' \emph{cascade-scoped revision}---``revision'' in the systems-approach sense, covering both verifier-detected repairs and user-requested refinements.

\subsection{Four-Layer Memory and the Session Lifecycle}
\label{sec:memory}

The mechanisms above guarantee consistency only within one session, while the lifecycle of \cref{sec:introduction} spans many sessions and the model is stateless---an LLM invocation sees only its own context. Without persistent state, every refinement session degenerates into regeneration from scratch, i.e., back to the first horn of the dilemma. The harness therefore organizes memory as a quadruple $M = (M_{\mathrm{w}}, M_{\mathrm{s}}, M_{\mathrm{e}}, M_{\mathrm{p}})$, partitioned by information lifetime: \emph{working memory} holds open-session state (enabling resume), \emph{semantic memory} holds stable knowledge (ADDIE, Bloom, rubric templates), \emph{episodic memory} holds session history and repair cases, and \emph{procedural memory} holds versioned policies such as the impact policy. Over one training program, the lifecycle is a chain of committed graph versions (superscript $G^{(v)}$ distinguishes versions from within-session steps $G_t$):
\[
G^{(1)} = \mathrm{init}(b), \qquad
G^{(v+1)} = \mathrm{commit}\Big(\mathrm{cascade}\big(\mathrm{apply}(G^{(v)}, r_v)\big)\Big),
\]
where $b$ is the initial design brief and $r_v$ the refinement request of session $v{+}1$. This chain is precisely what the four process metrics of RQ2 measure; an RQ1 benchmark run is the special case consisting only of $G^{(1)}$.

\subsection{Brief Elicitation}
\label{sec:elicitation}

The session chain starts from a design brief---the document that initiates a design project by defining the problem to solve, the goals, and the constraints~\cite{read2012designbrief}, corresponding in ISD terms to the output of needs assessment. In practice the brief does not exist at project start: the need takes shape through brainstorming between the training owner and stakeholders (project managers, mentors, the learners themselves), and designers often cannot articulate it before beginning. A system that only accepts finished briefs abandons its users at exactly the stage where they need support. Because brief formation is a multi-turn human-in-the-loop conversation---fundamentally different from a design run, which is a single complete pass over a settled brief---it is factored out as a module that precedes the design workflow. The elicitation agent maintains a conversation $h$; at each turn with user input $u$ it returns either the next clarifying question or a draft brief:
\[
\varepsilon(h, u) \in \{\, \mathrm{ask}(q), \; \mathrm{propose}(\hat{b}) \,\}.
\]
Only when the user approves the draft does $b := \hat{b}$ and the init session $G^{(1)} = \mathrm{init}(b)$ start the chain of \cref{sec:memory}; the conversation transcript is stored in working memory so the design session continues with the same context. Restricting the output to two actions (ask / propose) keeps the loop convergent and keeps the final decision with the designer.

\subsection{Agentic Retrieval}
\label{sec:agentic-rag}

The remaining requirement is domain knowledge: a generic agent cannot absorb the conventions of every domain---internal rubrics, exercise templates, code review standards---through a system prompt, yet hard-coding domain knowledge into prompts destroys the harness's generality. We resolve the tension by making domain knowledge \emph{pluggable}: each domain is packaged as a \emph{retrieval agent} with three deployer-supplied components---a corpus, a system prompt describing what the corpus answers well, and tools operating on that corpus. With registry $R = \{s_1, \dots, s_m\}$, extending to a new domain is just $R \leftarrow R \cup \{s_{m+1}\}$, changing no prompt or core harness code---the basis for the domain-transfer assumption discussed in \cref{sec:discussion}.

The design is \emph{agentic}~\cite{lewis2020rag,singh2025agenticrag} at two decision levels, both owned by agents rather than a fixed pipeline. At the usage level, retrieval is an on-demand tool of the specialist: no retrieval step runs at session start; while executing task $\tau$, the specialist $A_i$---the only party that knows what material is missing---makes the decision
\[
\delta_i(\tau, c) \;\in\; \{\varnothing\} \;\cup\; \big(2^R \times \{\text{single}, \text{seq}, \text{par}\} \times Q\big),
\]
i.e., either skip retrieval ($\varnothing$---retrieval costs time and tokens, so the default is no) or choose a set of retrieval agents, an execution strategy, and a query $q \in Q$ per agent; a default strategy is suggested by ADDIE task type (broad analytic/creative tasks favor parallel, narrowly scoped design tasks favor single), but the specialist decides. At the execution level, each selected retrieval agent runs its own bounded tool-use loop over its corpus to return exactly the requested material. The whole plan runs in a \emph{retrieval sandbox}: failures of one retrieval agent are isolated (they cannot crash the design session), tool-call budgets and result sizes are capped, and retrieval agents receive only a copy of the query---they never touch the graph or memory. The task's context pack becomes
\[
c_\tau = \rho(b) \;\cup\; \gamma(G_t) \;\cup\; \mu(M) \;\cup\; \mathrm{rag}\big(\delta_i(\tau, c)\big), \qquad \mathrm{rag}(\varnothing) = \varnothing,
\]
where $\rho(b)$ is context prepared from the brief, $\gamma(G_t)$ is graph-context retrieval, and $\mu(M)$ is memory retrieval---the latter two are always on and are distinct from external-corpus retrieval. Every decision $\delta_i$ and its outcome are logged in the trajectory, so the system's retrieval behavior is itself observable data in the case study.

\section{Case Study: An IT/Software-Engineering ISD Prototype}
\label{sec:case-study}

Benchmark evaluation (\cref{sec:experiments}) treats each scenario as an independent run---appropriate for comparing agents, but blind to the two characteristics established in \cref{sec:introduction}: real training programs are refined over many sessions (the core problem), and every domain carries its own documents and conventions (the domain-knowledge requirement). The case study fills that gap with a working application in our own enterprise context: internal IT/software-engineering training for new employees, interns, and junior engineers.

\paragraph{Application architecture and session lifecycle.}
The application wraps the core harness in a session-management layer that lets an alignment graph persist as a long-lived artifact at the level of a training program: \texttt{training\_program\_id} is the durable key of the graph, and an optional \texttt{session\_id} resumes an interrupted session. The first session (\textsl{init}) creates $G^{(1)}$; each subsequent session (\textsl{refinement}) loads the latest version, applies the refinement request, and commits $G^{(v+1)}$, exactly the version chain of \cref{sec:memory}. No new instrumentation is needed: trajectories, cascade events, repair history, and graph version diffs are already recorded by the harness.

\paragraph{Domain instantiation.}
Three representative programs were built: \emph{Java/OOP for junior engineers} (3 days, hybrid), \emph{software testing} (2 days, instructor-led), and a \emph{Git/agile mini-project} (2 days, for new project members). Each subject is packaged as a retrieval agent following the plug-and-play scheme of \cref{sec:agentic-rag}---a corpus (0--4 sample rubrics, assignment and starter-repository conventions, SOLID/code-review checklists, pull-request and sprint-simulation conventions), a system prompt describing the corpus scope, and a standard lookup toolset---and registered in $R$. The full workflow is demonstrated from the start of the chain: a vague request (``we need OOP training for juniors before they join the project'') passes through elicitation (\cref{sec:elicitation}) to an approved design brief; the init session generates $G^{(1)}$ with domain knowledge from the corpus; and refinement sessions then handle the requests that arise after deployment (add an objective, tighten a rubric, change the duration, switch the delivery format).

The case study also serves as the substrate of the RQ2 experiment: eight refinement requests were designed across the three programs, spanning the impact-policy spectrum from the narrowest target type (evaluation) to the broadest (needs context), so that the expected cascade scope varies rather than staying flat; the measurement protocol and the four control conditions are given in \cref{sec:experiments}.

\section{Experiments}
\label{sec:experiments}

\subsection{Experimental Setup}
\label{sec:setup}

\paragraph{RQ1: single-pass benchmark.}
We use ISD-Agent-Bench~\cite{isdagentbench}: synthetic scenarios spanning a Context~$\times$~ISD matrix, an evaluator that scores a 100-point ADDIE rubric (33 sub-steps $\rightarrow$ 13 items $\rightarrow$ 5 phases) plus a trajectory score, and multi-judge grading (GPT + Gemini) to reduce self-preference bias. AlignmentGraph-ISD is integrated as a seventh agent under the same input/output schema and compared with six agents: Baseline (single prompt), ReAct-ADDIE, ADDIE-Agent, Dick-Carey-Agent, RPISD-Agent, and EduPlanner. Two runs are executed on the same \texttt{test\_90} split---90 of 1{,}202 test scenarios, stratified over 10 domains $\times$ 3 difficulty levels $\times$ 3 scenarios per cell (seed 42)---with identical configurations except the agent model, covering two distinct points on the model-scale spectrum: \textbf{DeepSeek-v4-flash} (LLM) and \textbf{Qwen3.5-9B} (9B SLM). The judges are held fixed across both runs and use the benchmark's own default judge configuration, unmodified: \texttt{openai/gpt-4o-mini} and \texttt{google/gemini-2.5-flash-lite}. For fairness to the baselines, RQ1 runs the harness in its benchmark configuration: elicitation is not applicable (scenarios supply a fixed brief) and external-corpus retrieval is disabled (no agent accesses an external corpus). Each run has a manifest (model, judges, prompt versions, commit, seed, decoding configuration) for reproducibility.

\paragraph{RQ2: multi-session refinement experiment.}
The three case-study programs and eight refinement requests (\cref{sec:case-study}) run under the four conditions of \cref{tab:rq2-conditions}. All four conditions receive the same information (the previous design plus the refinement request); the only difference is the mechanism. Unlike RQ1, the two harness conditions here run the full architecture, including agentic retrieval over the IT/software-engineering corpora; elicitation is demonstrated in the case study (\cref{sec:case-study}, results in \cref{sec:case-study-results}).

\begin{table*}[t]
  \caption{The four control conditions of the RQ2 experiment.}
  \label{tab:rq2-conditions}
  \begin{tabular}{@{}p{0.22\textwidth}p{0.19\textwidth}p{0.51\textwidth}@{}}
    \toprule
    Condition & Agent & Refinement mechanism \\
    \midrule
    \texttt{harness\_targeted\_refine} & AlignmentGraph-ISD & refinement request $\rightarrow$ impact-policy cascade $\rightarrow$ targeted repair (proposed) \\
    \texttt{harness\_full\_regen} & AlignmentGraph-ISD & same information, all nodes invalidated $\rightarrow$ full regeneration (isolates the scoping mechanism) \\
    \texttt{react\_full\_regen} & ReAct-ADDIE (strongest RQ1 baseline) & full re-run from the updated scenario --- horn (1) of the dilemma \\
    \texttt{react\_targeted\_refine} & ReAct-ADDIE + \texttt{refine} endpoint & prompt with old scenario + old output + change request, no alignment control --- horn (2) of the dilemma \\
    \bottomrule
  \end{tabular}
\end{table*}

Four \emph{process metrics} are measured with shared output-level tooling so they compare across agents: (1)~\emph{cascade scope}---the fraction of nodes marked as affected (meaningful for the two harness conditions); (2)~\emph{edit locality}---the fraction of ADDIE output subsections that change between consecutive rounds, measured by content hashing (all conditions); (3)~\emph{cost per refinement round}---tokens and latency, absolute and relative to the init session of the same chain; and (4)~\emph{alignment preservation}---the evaluator's alignment rubric item (D1) and the ADDIE score before and after each round, graded by the same judges for all four conditions; the two harness conditions additionally report internal audit coverage (\textsl{measures}/\textsl{supports}/\textsl{evaluates}) as an appendix. The model is the same SLM as in RQ1 (Qwen3.5-9B). Every session records a trajectory and a manifest.

\subsection{RQ1 Results}
\label{sec:rq1-results}

\Cref{tab:rq1-results} summarizes the two-judge mean scores of the seven agents on \texttt{test\_90} for both model runs.

\begin{table*}[t]
  \caption{ISD-Agent-Bench results on \texttt{test\_90} (two-judge means; Total = composite, ADDIE = artifact rubric, Traj.\ = trajectory; \textbf{bold} = best in column).}
  \label{tab:rq1-results}
  \begin{tabular}{@{}lrrrrrr@{}}
    \toprule
    & \multicolumn{3}{c}{DeepSeek-v4-flash (LLM)} & \multicolumn{3}{c}{Qwen3.5-9B (SLM)} \\
    \cmidrule(lr){2-4} \cmidrule(lr){5-7}
    Agent & Total & ADDIE & Traj. & Total & ADDIE & Traj. \\
    \midrule
    Baseline          & 82.15 & 83.93 & 77.99 & 78.47 & 75.02 & 86.52 \\
    EduPlanner        & 77.49 & 76.45 & 79.92 & 54.51 & 42.03 & 83.65 \\
    ADDIE-Agent       & 80.16 & 79.44 & 81.83 & 80.17 & 76.63 & 88.43 \\
    RPISD-Agent       & 79.09 & 79.83 & 77.38 & 78.77 & 77.16 & 82.54 \\
    Dick-Carey-Agent  & 82.70 & 82.41 & 83.37 & 82.03 & 79.56 & 87.80 \\
    ReAct-ADDIE       & 81.25 & 78.67 & \textbf{87.26} & 79.14 & 73.77 & \textbf{91.66} \\
    \textbf{AlignmentGraph-ISD} (ours) & \textbf{84.07} & \textbf{85.29} & 81.20 & \textbf{84.65} & \textbf{84.02} & 86.11 \\
    \bottomrule
  \end{tabular}
\end{table*}

\paragraph{Overall position.}
AlignmentGraph-ISD leads both the composite and the ADDIE score in both runs: 84.07 with DeepSeek (runner-up Dick-Carey-Agent, 82.70) and 84.65 with Qwen (runner-up Dick-Carey-Agent, 82.03). On Qwen it leads under both judges (GPT 85.05, Gemini 84.25); on DeepSeek it leads under the GPT judge (86.62), while the Gemini judge ranks ReAct-ADDIE higher by 0.30 points. It also leads at all three difficulty levels in both runs and achieves a 100\% success rate in both (DeepSeek run: Dick-Carey-Agent 88.9\%, EduPlanner 94.4\%, Baseline 97.8\%; Qwen run: Baseline 98.9\%, EduPlanner 97.8\%).

\paragraph{Statistical tests.}
Wilcoxon signed-rank tests paired by scenario on the composite score (pairs with failed runs excluded), with Holm correction over the six comparisons per run, show AlignmentGraph-ISD above all six baselines in both runs with $p_{\mathrm{Holm}} < 0.01$ for all 12 comparisons. Mean paired differences (bootstrap 95\% CIs excluding zero in every comparison) range from $+1.53$ (vs.\ Dick-Carey-Agent) to $+6.49$ (vs.\ EduPlanner) on DeepSeek, and from $+2.62$ (Dick-Carey-Agent) to $+6.16$ (Baseline) on Qwen, apart from EduPlanner's collapse on the SLM ($+30.11$).

\paragraph{SLM effect.}
Moving from the LLM to the SLM leaves the harness essentially unchanged (84.07 $\rightarrow$ 84.65), while the single-prompt Baseline loses 3.68 points (82.15 $\rightarrow$ 78.47) and EduPlanner loses 22.98; the paired margin over Baseline accordingly widens from $+1.86$ to $+6.16$. A per-phase decomposition shows where the difference arises: on Qwen, Baseline degrades most in the long-form content phases---Development 66.2 and Implementation 61.3---while the harness holds 85.1 and 82.2, leading four of five phases (Design, Development, Implementation, Evaluation; Analysis trails RPISD-Agent by 0.9). This pattern is consistent with the harness-over-model argument~\cite{cho2026harness}: verification--repair and graph-driven coordination directly compensate for missing model capability.

\paragraph{Consistent weakness.}
The trajectory score is below ReAct-ADDIE in both runs (81.20 vs.\ 87.26 on DeepSeek; 86.11 vs.\ 91.66 on Qwen); we return to this in \cref{sec:discussion}.

\subsection{RQ2 Results}
\label{sec:rq2-results}

[RQ2 RESULTS PENDING --- table of 4 conditions $\times$ 4 metrics $\times$ 8 rounds. Mechanism-level expectation validated in a dress rehearsal: cascade scope of \texttt{harness\_targeted\_refine} stratifies from 0.17 to 1.0 by target node type while both full-regeneration conditions stay $\approx$ 1.0; the narrowest round takes 2 LLM calls versus 6 for full regeneration.]

Mechanism-level evidence is already available from RQ1: cascade-scoped revision (verifier detection $\rightarrow$ scoping $\rightarrow$ regeneration) operated inside every benchmark run, is recorded in the trajectories, and is reflected in the repair success rate and the 100\% success rate in both \texttt{test\_90} runs.

\subsection{Case-Study Results}
\label{sec:case-study-results}

Qualitative demonstration of the full workflow described in \cref{sec:case-study} (elicitation $\rightarrow$ init $\rightarrow$ refinement sessions). [TRANSCRIPT PENDING --- one representative session chain to be selected after the production runs.]

\section{Discussion}
\label{sec:discussion}

\paragraph{Interpretation.}
RQ1 shows that a refinement-oriented harness does not trade away first-pass quality---on the contrary, it leads both the composite and the ADDIE score in both runs, and is nearly invariant to swapping the LLM for a 9B SLM (84.07 $\rightarrow$ 84.65). The trajectory-score gap to ReAct-ADDIE is consistent (81.20 vs.\ 87.26 on DeepSeek; 86.11 vs.\ 91.66 on Qwen) and reflects a current limitation of our logging and dynamic orchestration: the trace records the outcome of each step more than the decision process behind it. [RQ2 INTERPRETATION PENDING --- to be written from the four-condition table.]

\paragraph{Threats to validity.}
\emph{Internal:} prompts, models, and judges all influence scores; we fix manifests, use identical decoding configurations, grade with multiple judges, and randomize grading order. The \texttt{react\_targeted\_refine} baseline uses a refinement prompt written by us; we mitigate this by publishing the prompt and designing it as favorably as possible for the baseline (the model may return \texttt{unchanged} per phase to save cost). The ReAct baseline in RQ2 uses an English translation of the benchmark's original Korean prompt (removing the language confound; both versions are published with semantics preserved).
\emph{Construct:} the benchmark and the process metrics measure artifact and process quality, not learner outcomes; our conclusions are bounded accordingly.
\emph{External:} benchmark scenarios are synthetic and multi-domain, while the case study covers a single domain (IT/software engineering); transferring domains by swapping corpora, templates, and verifier policies is an assumption that requires further validation.
\emph{Cost:} multi-agent generation plus repair costs more than a single prompt on the first pass; we report the full quality/cost trade-off rather than quality alone.

\paragraph{Broader implications.}
The pattern ``typed dependency graph + verifier + cascade-scoped revision'' is not education-specific: any setting that generates many interdependent artifacts (technical documentation, process definitions, curricula) faces the same refinement dilemma. The SLM result---the harness's score essentially unchanged while the single-prompt baseline drops 3.68 points and the paired margin widens from $+1.86$ to $+6.16$---strengthens the harness-over-model argument for low-cost enterprise deployment~\cite{cho2026harness,belcak2025slm}.

\section{Conclusion}
\label{sec:conclusion}

The real lifecycle of a corporate training program---a brief that emerges from stakeholder brainstorming, followed by continuous post-deployment refinement---poses a dilemma that single-pass LLM agents cannot resolve, and its root cause is the absence of an explicit representation of alignment. AlignmentGraph-ISD makes alignment the first-class graph state of an agent harness and organizes every module---elicitation, orchestration, retrieval, verification, and multi-session memory---around that state, so that a refinement request is scoped, regenerated exactly where needed, re-verified, and committed as a new version (cascade-scoped revision). Evidence on ISD-Agent-Bench confirms leading first-pass quality under both an LLM and a 9B SLM; the four-condition experiment [RQ2 RESULTS PENDING] quantifies both horns of the dilemma directly. Future work includes enriching the trajectory to close the trajectory-score gap, validating the domain neutrality of the alignment graph, and evaluating elicitation with practicing designers.


\section*{Data Availability}
ISD-Agent-Bench, including all six baseline agents and the evaluation
pipeline, is publicly available~\cite{isdagentbench}. Run manifests
(model, judge, prompt versions, seeds, and decoding configurations),
the baseline refinement prompts used in RQ2, and the case-study design
briefs and refinement requests will be released with the camera-ready
version of this paper.

\bibliographystyle{ACM-Reference-Format}
\bibliography{references}

\end{document}
